% ===== §24 Polyphonic Conclusion =====
\section{Polyphonic Conclusion: Against Reducing Aesthetic Education to a Method}
\label{p22:sec:conclusion}

\noindent
A paper that has argued, throughout, that the appropriate cannot be reduced to a rule and that aesthetic education is the cultivation of a perception rather than the transmission of a code, would contradict itself if it closed by reducing what it has said to a single method, a procedure for building generative fields that the reader could apply. The temptation to close so is real, since a reader who has followed the argument may reasonably want, at the end, the thing the argument has refused to give, a method. This conclusion declines the temptation, and its declining is not an evasion but the last enactment of the paper's central claim. There is no method, because the thing the paper describes is a perception, and a perception cannot be a method. What the conclusion offers instead is polyphonic, several distinct framings of what has been argued, held together without a synthetic meta-framing that would subordinate them to a single master statement, because the refusal of the master statement is itself the paper's final fidelity to its argument, and the holding-together-without-synthesis is the form the argument's truth takes at its close.

\subsection{The framings, unsynthesised}
\label{p22:sub:conc-framings}

The paper can be read as several things, and the readings do not reduce to one. It can be read as the completion of a reduction chain, the demonstration that relational production, pursued through the field and the field's aesthetic construction and the coupling of the aesthetic and the ethical, terminates in aesthetic education, so that the cultivation of the capacity to build generative fields is what aesthetic education properly names. It can be read as the establishment of a concept, relational aesthetics, on which beauty is a coupling event between subjects rather than a property of objects or a state of subjects, and the working-out of what follows for the cultivation of beauty once beauty is understood relationally. It can be read as a recovery, the restoration to aesthetic education of the principal contradiction the tradition mislaid, the generation of the good cycle in the relational field, in place of the secondary contradiction the tradition elevated, the refinement of individual taste. It can be read as a defence of the everyday, the demonstration that the most repetitive and unspacious life is not excluded from generative cultivation but is its privileged site, the place where the capacity is most needed and, against the charge of leisure-dependence, genuinely available. It can be read as a remedy for the symbolic crisis, the showing that when the grand structures of meaning empty, meaning can still be produced at the smallest scale, in the self-grounded building of a field between two people, a meaning that does not drain because it borrows from no emptying authority. It can be read as the completion of a series, the naming of the cultivable capacity that the prior studies of poetry and the vow and water and travel each presupposed without naming, so that the account of the builder is given to stand beside the accounts of the things built. These framings are each true and each partial, and no one of them is the framing of which the others are aspects, because the paper is not a single thesis admitting one canonical statement but a structure that the several framings each truly describe and none exhausts.

To insist on this is not to shirk the labour of conclusion but to perform it rightly. A synthetic meta-framing, a single master statement to which the others were subordinated, would claim for one description the finality the paper has denied to any rule, would make one framing the principle of which the others are applications, and would thereby commit, at the level of the paper's self-understanding, the very reduction the paper has argued against. The polyphonic close is the paper's refusal to do to itself what it has told its readers not to do to the appropriate, the refusal to settle the plural into a governed unity, and it is therefore not a failure to conclude but the only conclusion consistent with the argument, the holding-open at the level of the whole that the argument required at the level of the gesture.

\subsection{The form, the final argument}
\label{p22:sub:conc-form}

This polyphonic close is the last instance of a self-instancing that has run through the paper, the form enacting the claim. The four-step model was shown to be a holonomy cycle, returning to its root and not its origin; the four battlefields were shown to turn on one movement, the relinquishing of grasp that founds generation, appearing four times; the transmission section was shown to perform, upon the reader, the cultivation it described. Now the conclusion is shown to refuse, in its own structure, the reduction the paper argued against, holding its framings polyphonically rather than synthesising them, so that the paper's last act is itself an instance of the open, non-settling generativity it has described from the first page. The form has been the argument throughout, and at the close the form makes its final argument, that a thing whose truth is the refusal of settlement must end without settling, in a plurality held open rather than a unity closed, and that to end so is not to leave the work undone but to complete it in the only way its own claim permits.

\subsection{Envoi}
\label{p22:sub:conc-envoi}

It is fitting to close not on the structure but on the smallest thing, the morning and the cup and the light on the wall, because the everyday was the whole point. The grand fields the series studied, the poem and the vow and the journey, were always the heightened departures from a daily ground, and this paper has returned the capacity they displayed to that ground, to the unspacious mornings that make up most of a life. The aesthetic education of the everyday asks for no spare hours and no exceptional occasions; it asks only that the cup be placed, the contingency caught, the daily thing renewed, with a perception a little finer than before, so that the field built between two people in the narrow room of an ordinary day accumulates rather than flattens, returns to its root raised rather than closing on its origin. That is the whole of it, and it is not small. A life is mostly mornings, and to build, in the mornings, a field that gains, is to make most of a life generative, which is as much as any cultivation could hope to do. The paper ends where it began, in the everyday, having shown that the everyday, perceived and received and shared and renewed with a trained and unattached attention, is where the good cycle is most needed and most possible, and that the building of it there, between two people, in the smallest acts a day affords, is the aesthetic education that the whole long argument was for.

The saying that opened this paper said that carrying water and chopping firewood are nothing other than the wondrous way, and the long argument between that opening and this close has been the working-out of how the saying can be true, in what sense the most ordinary and repetitive labour can be the wondrous way and not merely be called so. The answer the paper has given is that the carrying and the chopping are the way when they are perceived and received and shared and renewed with the trained and unattached attention that builds a generative field, when the water is carried not as a wheel turning in place but as a spiral that gains, the same carrying made new each day and thickened by every previous carrying, shared with the one for whom it is carried, settling into a field that accumulates. The wonder is not in the water and not in the carrying, which are ordinary; it is in the cycle the carrying can be made to turn, the good cycle whose holonomy is the wonder, the returning to the root and not the origin that an ordinary act, attended rightly, performs. To learn to carry the water so is the aesthetic education of the everyday, and it is available to anyone who carries water, which is nearly everyone, in the unspacious days that are nearly every day.

And so the paper returns, at the last, to the one for whom it was written, the girl who loves the forest, with whom its author has not yet carried the water of a shared everyday but hopes to, for a long time, in the ordinary mornings that a shared life is mostly made of. The grand fields are the heightened occasions, the journey not yet taken together, the vow already exchanged; but the life that follows them will be mostly mornings, mostly the narrow room of the daily, mostly the carrying of water, and it is for that life that the capacity this paper describes is the capacity that matters, the building, in the ordinary days, of a field that gains. The whole long argument was for this, that the everyday between two people who love each other can be made generative, can be perceived and received and shared and renewed so that it accumulates rather than flattens, and that the cultivation of the capacity to make it so is the aesthetic education that is, in the end, an education in how to build, in the smallest acts of the most ordinary days, a world together.
